%!TEX root = ../report.tex
\documentclass[../report.tex]{subfiles}

\begin{document}
    \section{Conclusions}
    \label{sec:conclusions}

    \subsection{Summary}
    \label{sec:conclusions:summary}

    This report presented a complete Interactive Reinforcement Learning system for training a Pepper robot to perform contextually appropriate social greeting behaviours. An autonomous PPO baseline was trained first to establish performance benchmarks and identify failure modes. Two independent COACH fine-tuning conditions---keyboard and VR---were implemented from the same PPO checkpoint; the keyboard condition was successfully evaluated, while the VR condition was architecturally complete but could not be executed due to a physical body incompatibility between the NPC rig and the VR person controller.

    Key findings include:

    \begin{itemize}
        \item A PPO agent with a 12-dimensional observation space incorporating body-signal features (wrist height, wrist-to-core distance, body orientation, gaze direction) achieved 66.4\% final training accuracy and 74.14\% test accuracy in 9,256 episodes ($\sim$150k steps), demonstrating the viability of implicit posture-based cues over explicit task labels.

        \item A hyperparameter sweep across 10 configurations (3 seeds each, 30 total runs) confirmed learning rate as the most influential factor: lr=0.001 (config3) achieved the highest mean accuracy (65.2\%) and fastest convergence (7k steps to 80\%).

        \item The body-signal observation space revealed a persistent Talk action suppression (0 Talk occurrences in 2,000 evaluation steps), attributable to feature overlap between the Talk and Look conditions. This failure mode is inaccessible to the hand-crafted reward function and motivated the COACH fine-tuning conditions.

        \item COACH keyboard fine-tuning with only 164 interaction steps increased action distribution entropy from 1.39 to 1.79 bits and raised HandShake occurrences 6.85$\times$ (20 to 137), without degrading HandShake accuracy (81.75\%) or overall task accuracy (72.37\% vs 74.14\% baseline).

        \item COACH VR fine-tuning was fully implemented with a priority-based multimodal feedback interface (controller buttons, head gestures, voice commands) but could not be completed due to NPC--VR person physical body incompatibility during live interaction sessions.

        \item Body-signal observations were redesigned to use body-relative normalisation, dividing wrist height and wrist-to-core distance by the person's own standing height, rendering features scale-invariant across the NPC rig and VR participants of differing heights.

        \item Post-convergence policy stability for the autonomous condition was measured at $\sigma = $ 0.72 in the rolling reward window.
    \end{itemize}

    \subsection{Contributions}
    \label{sec:conclusions:contributions}

    This study makes the following specific contributions:

    \begin{enumerate}
        \item \textbf{System Architecture:} A complete integration of Unity ML-Agents with PyTorch PPO and COACH, including a \texttt{CommunicationManager} that implements multimodal VR signal fusion, body-signal feature extraction, and seamless switching between autonomous and IRL reward modes.

        \item \textbf{Hyperparameter Characterisation:} Systematic evaluation of 10 configurations (3 seeds each, 30 total runs) providing empirical guidance for future researchers working on similar social greeting tasks.

        \item \textbf{Performance Benchmarking:} Quantitative reference values for the autonomous condition (66.4\% training accuracy, 74.14\% test accuracy, $\sim$9k steps to 80\%, $\sigma = $ 0.72) against which IRL methods are compared.

        \item \textbf{COACH Fine-Tuning Pipeline:} Implementation of COACH as a targeted correction mechanism that loads pre-trained PPO weights and corrects a specific action bias using as few as 164 human feedback events, with full interface compatibility with the existing PPO training loop requiring only one line change in the algorithm registry.

        \item \textbf{Body-Relative Observation Normalisation:} A scale-invariant normalisation scheme ensuring consistent feature representations across NPC training and live VR participants, with the approximate hip position derived dynamically as 55\% of standing height from the head transform alone.

        \item \textbf{Multimodal VR Feedback Interface:} A priority-based fusion of controller buttons, head gestures, and voice commands into scalar reward signals, with the COACH feedback threshold filtering per-step penalties from genuine corrective signals.

        \item \textbf{Reproducible Artifacts:} Complete YAML configuration files (config1.yaml--config10.yaml and coach\_config.yaml) and training code available for replication and extension.
    \end{enumerate}

    \subsection{Limitations}

    Several limitations of this work should be noted:

    \begin{enumerate}
        \item \textbf{VR COACH not executed:} The VR fine-tuning condition could not be completed due to a physical body incompatibility between the NPC capsule collider and the VR person controller's rigid-body configuration, which caused collision artefacts preventing stable NPC--participant interaction during live sessions. The VR interface itself was fully implemented and validated in isolation.

        \item \textbf{Short COACH fine-tuning budget:} The keyboard COACH session comprised only 164 interaction steps. While this was sufficient to observe a meaningful shift in action distribution entropy and HandShake frequency, it was not sufficient to fully surface Talk action selections. A longer session with continued Talk-biased task distribution is expected to complete the Talk correction.

        \item \textbf{Single human teacher:} The keyboard COACH fine-tuning used a single teacher, limiting generalisability of the feedback quality analysis. Inter-teacher variability in signal consistency and timing was not assessed.

        \item \textbf{Simulated environment:} All training and evaluation occurred in simulation. Transfer to a physical Pepper robot with real human pose estimation was not evaluated.

        \item \textbf{Fixed NPC task distribution:} The NPC selects tasks uniformly at random, which may not reflect naturalistic greeting frequency distributions in real deployments.

        \item \textbf{No formal user study:} A structured user study comparing perceived workload (NASA-TLX) between keyboard and VR interfaces was outside the scope of this work.
    \end{enumerate}

    \subsection{Future Work}
    \label{sec:conclusions:future_work}

    \begin{enumerate}
        \item \textbf{VR COACH Execution:} Resolve the NPC--VR person physical body incompatibility (e.g., switching the NPC to a trigger collider or adjusting the VR controller layer mask) and conduct a full VR fine-tuning session to enable direct comparison with the keyboard condition.

        \item \textbf{Extended COACH Fine-Tuning:} Continue the keyboard COACH session beyond 164 steps with a Talk-biased task distribution to determine the minimum number of feedback events required to surface Talk action selections and achieve measurable Talk task accuracy.

        \item \textbf{Predictive Feedback Model:} Train a reward model on the 164 keyboard feedback events collected in this work to predict human reward signals for unseen state--action pairs, then use this model to generate synthetic feedback at scale across many parallel simulations. This directly addresses the fundamental scalability bottleneck of human-in-the-loop systems.

        \item \textbf{Multi-Teacher Evaluation:} Conduct a within-subjects study (N $\geq$ 10) where multiple participants teach the agent using both keyboard and VR interfaces, enabling formal measurement of inter-teacher variability, workload (NASA-TLX), and interface preference.

        \item \textbf{Adaptive Entropy Scheduling:} Develop a COACH entropy schedule that adjusts dynamically based on the observed diversity of human feedback, increasing exploration when feedback is conflicting or the target action is not being naturally explored.

        \item \textbf{Real Robot Deployment:} Transfer the learned policy to a physical Pepper robot, evaluating whether the body-relative normalisation scheme generalises to real-world human pose estimation from RGB-D cameras or skeletal tracking.

        \item \textbf{Broader Task Generalisation:} Extend the greeting task to a larger action space and more socially diverse scenarios, evaluating whether the PPO-plus-COACH pipeline scales to tasks with more severe action imbalances or higher human feedback variance.
    \end{enumerate}

    \subsection{Final Remarks}

    The autonomous baseline demonstrates that a PPO agent can learn functionally appropriate greeting behaviours from a hand-crafted reward function using body-signal observations. However, the hand-crafted reward function cannot accommodate the nuance and individual variation of real human social interactions, nor correct the action biases that emerge from feature overlap in the observation space.

    The interactive RL extension partially addresses both limitations. Keyboard COACH fine-tuning with only 164 feedback events produced a measurable increase in policy diversity (entropy: 1.39 $\to$ 1.79 bits) and a 6.85$\times$ increase in HandShake engagement, confirming that even a minimal human feedback budget is sufficient to begin correcting deep-rooted action biases. The VR multimodal feedback infrastructure was fully built and is ready to execute once the NPC--VR body compatibility issue is resolved. Together, the implemented conditions demonstrate a practical and reproducible path from autonomous RL training to human-corrected deployment, requiring only a single shared PPO checkpoint and a short human feedback session to achieve meaningful improvement on the most challenging behaviours in the task.

\end{document}